\chapter{Объектно-ориентированное программирование}
\label{ch:03}

Предыдущая глава дала строительные блоки языка: классы и интерфейсы, массивы,
строки, записи и перечисления. Но программа, выполняющая инструкции сверху
вниз, не умеет ни выбирать, ни повторять, а набор классов сам по себе
ещё не система. Поэтому глава начинается с ветвлений и циклов, а затем
переходит к принципам объектно-ориентированного программирования: вы научитесь
строить иерархии классов, отличать отношение «является» от отношения
«содержит», прятать состояние объекта за методами и писать код, одинаково
работающий с любым потомком.

\section{Ветвление}

Любая полезная программа принимает решения. Навигатор смотрит на дорогу:
перекрыта~--- ведёт в объезд, свободна~--- прямо. Развилку в коде описывает
оператор \texttt{if}, а когда вариантов больше двух, ветки выстраивают в
цепочку из \texttt{if}, нескольких \texttt{else if} и \texttt{else}.

Условие в скобках обязано иметь тип \texttt{boolean}: ни число, ни ссылка
условием быть не могут, а опечатка \texttt{if (x = 5)} даст ошибку компиляции.
Ветки проверяются сверху вниз, и выполняется первая подошедшая, поэтому порядок
важен: если проверка \texttt{score >= 75} стоит в цепочке раньше
\texttt{score >= 90}, отличник получит «хорошо».

\begin{oshibka}
\textbf{Частая ошибка: пропущенные фигурные скобки.} Без блока \texttt{\{\}}
оператор \texttt{if} управляет ровно одним следующим оператором, а не всеми
строками с тем же отступом: в коде
\texttt{if (x > 5) System.out.print("A"); System.out.print("B");} буква
\texttt{B} печатается при любом \texttt{x}. Отступы для компилятора ничего не
значат~--- ставьте скобки даже вокруг одной строки.
\end{oshibka}

Когда ветвление выбирает не действие, а значение, короче \term{тернарный
оператор} \texttt{?:}~--- строка \texttt{int max = (a > b) ? a : b;} заменяет
собой \texttt{if} из четырёх строк с двумя присваиваниями.

\subsection{Оператор switch}

Цепочка из семи \texttt{else if}, сравнивающих одну переменную с разными
значениями, читается плохо; для этого случая есть \texttt{switch}, классическая
форма которого (листинг~\ref{lst:03-1}) требует аккуратности.

\begin{listing}[H]
\begin{minted}{java}
switch (day) {
    case 1:
        System.out.println("Понедельник");
        break;
    case 2:
    case 3:
        System.out.println("Вторник или среда");
        break;
    default:
        System.out.println("Другой день");
}
\end{minted}
\caption{Классический switch и проваливание}
\label{lst:03-1}
\end{listing}

Слово \texttt{break} завершает работу оператора. Забыли~--- и выполнение
\term{проваливается} (\eng{fall-through}) в следующую ветку и идёт дальше до
первого \texttt{break}; намеренно так поступают редко~--- разве что объединяя
варианты двумя метками \texttt{case} подряд. Java~14 убрала ловушку: в
стрелочной форме проваливания нет, а сам \texttt{switch} работает и как
оператор, и как выражение (листинг~\ref{lst:03-2}).

\begin{listing}[H]
\begin{minted}{java}
switch (level) {           // как оператор: break не нужен
    case 1 -> System.out.println("Начинающий");
    case 2, 3 -> System.out.println("Средний уровень");
    default -> System.out.println("Неизвестный уровень");
}

String name = switch (level) {      // как выражение
    case 1 -> "один";
    default -> { yield "другое"; }  // yield - итог блока
};
\end{minted}
\caption{Стрелочный switch: оператор и выражение}
\label{lst:03-2}
\end{listing}

Одиночное выражение после стрелки становится результатом само; если ветка~---
блок, результат называют явно словом \texttt{yield}.
\texttt{switch}-выражение обязано покрывать все случаи, иначе компилятор
потребует ветку \texttt{default}. Java~21 разрешила ветвиться не по значению, а
по типу и добавила к образцу условие \texttt{when} (листинг~\ref{lst:03-3}).

\begin{listing}[H]
\begin{minted}{java}
static String describe(Object obj) {
    return switch (obj) {
        case null -> "ничего нет";
        case Integer i when i > 100 -> "большое: " + i;
        case Integer i -> "число: " + i;
        case String s -> "строка длиной " + s.length();
        default -> "неизвестный тип";
    };
}
\end{minted}
\caption{switch с сопоставлением по образцу}
\label{lst:03-3}
\end{listing}

Запись \texttt{case Integer i} совмещает проверку типа и объявление уже
приведённой переменной. Порядок веток значим: образец с условием \texttt{when}
обязан стоять раньше такого же образца без условия. Принимает \texttt{switch}
целые типы, \texttt{String}, перечисления и~--- при сопоставлении с
образцом~--- любой объект; тип \texttt{long} недопустим.

\section{Циклы}

Если ветвление~--- развилка на дороге, то цикл~--- круговое движение: действия
повторяются, пока не выполнится условие выхода. Форм цикла четыре, и в
листинге~\ref{lst:03-4} каждая печатает числа 0, 1, 2.

\begin{listing}[H]
\begin{minted}{java}
int i = 0;
while (i < 3) {          // условие проверяется ДО тела
    System.out.println(i++);
}
int j = 0;
do {                     // тело выполнится хотя бы раз
    System.out.println(j++);
} while (j < 3);
for (int k = 0; k < 3; k++) {  // управление в заголовке
    System.out.println(k);
}
for (int n : new int[]{0, 1, 2}) {  // for-each: без индекса
    System.out.println(n);
}
\end{minted}
\caption{Четыре формы цикла}
\label{lst:03-4}
\end{listing}

Цикл \texttt{while} проверяет условие перед телом, поэтому тело может не
выполниться ни разу; \texttt{do-while} проверяет после, и одно выполнение
гарантировано~--- так удобно запрашивать ввод. У \texttt{for} всё управление в
заголовке: инициализация выполняется один раз, условие~--- перед каждым
проходом, шаг~--- после тела. Форма \texttt{for-each} читается лучше всех, но
индекс в ней недоступен.

Оборвать цикл досрочно позволяет \texttt{break}, пропустить проход~---
\texttt{continue}. Чтобы адресовать их не своему, а внешнему циклу, перед
циклом ставят \term{метку} (листинг~\ref{lst:03-5}).

\begin{listing}[H]
\begin{minted}{java}
search:                          // метка внешнего цикла
for (int r = 0; r < matrix.length; r++) {
    for (int c = 0; c < matrix[r].length; c++) {
        if (matrix[r][c] == target) {
            row = r;
            col = c;
            break search;        // выход из ОБОИХ циклов
        }
    }
}
\end{minted}
\caption{break с меткой}
\label{lst:03-5}
\end{listing}

Без метки \texttt{break} прервал бы только внутренний цикл, внешний продолжил
бы перебирать строки, и пришлось бы заводить флаг. Метка работает и с
\texttt{continue}: запись \texttt{continue search;} из циклов не выходит, а
начинает следующий проход внешнего цикла, пропустив остаток внутреннего. Но
метка~--- не переход в произвольную точку программы: оператора \texttt{goto} в
Java нет.

\begin{oshibka}
\textbf{Частая ошибка: удаление элементов внутри \texttt{for-each}.} Вызов
\texttt{list.remove(x)} прямо в цикле даёт исключение
\texttt{ConcurrentModificationException}: скрытый итератор обнаруживает, что
коллекцию изменили в обход него. Удаляйте явным итератором через
\texttt{it.remove()} или вызовом \texttt{list.removeIf(...)}.
\end{oshibka}

\section{Наследование}

Объектно-ориентированное программирование стоит на четырёх принципах:
\term{абстракция}~--- выделение существенного и сокрытие второстепенного;
\term{наследование}~--- иерархия, в которой потомок получает всё, что есть у
родителя; \term{инкапсуляция}~--- сокрытие состояния за методами;
\term{полиморфизм}~--- единый вызов, дающий разное поведение. Абстракцию мы
применяли, объявляя абстрактные классы; остальные три разберём по очереди.

В природе все животные дышат, едят и двигаются, а виды лишь уточняют общие
черты. Так же и в коде: общее поведение описывают в родительском классе,
потомок добавляет своё. Потомка объявляют словом \texttt{extends}, и расширить
он может ровно один класс~--- в Java принято \term{одиночное наследование}
классов, зато интерфейсов допустимо реализовать сколько угодно. У корня любой
иерархии (рис.~\ref{fig:03-1}) стоит класс \texttt{Object}.

\begin{figure}[htbp]\centering
\begin{forest} hierarchy
[Object
  [Animal
    [Dog]
    [Cat]
    [{Bird (abstract)}
      [Owl]
      [Penguin]]]
]
\end{forest}
\caption{Иерархия классов: потомки уточняют поведение предка}\label{fig:03-1}
\end{figure}

Как связаны родитель и потомок в коде, показывает листинг~\ref{lst:03-6}:
смотрите на \texttt{this.name}, \texttt{super(name)} и \texttt{super.speak()}.

\begin{listing}[H]
\begin{minted}{java}
class Animal {
    protected String name;

    public Animal(String name) {
        this.name = name;    // слева поле, справа параметр
    }

    public void speak() {
        System.out.println(name + " издаёт звук");
    }
}

class Dog extends Animal {
    public Dog(String name) {
        super(name);         // конструктор родителя
    }

    @Override
    public void speak() {
        super.speak();       // сначала поведение родителя
        System.out.println(name + " лает");
    }
}
\end{minted}
\caption{Родительский класс и его потомок}
\label{lst:03-6}
\end{listing}

Слово \texttt{this} означает «текущий объект»: запись \texttt{this.name}
различает поле и одноимённый параметр. Слово \texttt{super}~--- «родительская
часть»: \texttt{super(...)} вызывает конструктор родителя,
\texttt{super.speak()}~--- его версию переопределённого метода.

Отсюда следует порядок создания объекта: конструктор потомка первым делом
вызывает конструктор родителя, тот~--- своего, и так до \texttt{Object}, а тела
выполняются в обратном порядке~--- сначала \texttt{Object}, потом
\texttt{Animal}, потом \texttt{Dog}. Вызов \texttt{super(...)} обязан быть
первой инструкцией (Java~25 разрешает предварять его вычислениями, не
обращающимися к \texttt{this}); не написали~--- компилятор вставит
\texttt{super()} без аргументов, и если у родителя такого конструктора нет,
будет ошибка компиляции.
Той же первой инструкцией может быть \texttt{this(...)}~--- вызов другого
конструктора этого же класса.

\section{Композиция, агрегация и зависимость}

Наследование выражает отношение \term{is-a} («является»): собака~--- это
животное. Но чаще один класс \emph{содержит} другой или просто
\emph{пользуется} им.

Представьте квартиру. Комнаты~--- её неотъемлемая часть: снесут дом~--- исчезнут
и комнаты. Это \term{композиция}. Мебель~--- другое дело: диван можно вынести и
поставить в другую квартиру, он существует сам по себе. Это \term{агрегация}. А
курьер, заглянувший на пять минут, частью квартиры не становится~--- это
\term{зависимость}, связь «использует», а не «содержит». В коде класса
\texttt{Car} эти три отношения выглядят так: двигатель~--- поле, созданное
внутри самим классом (\texttt{private final Engine engine = new Engine();});
пассажиры~--- поле, полученное готовым в конструкторе; механик~--- лишь
параметр метода \texttt{repair}: ссылку на него класс не хранит.

\subsection{Одна задача~--- два решения}

Кажется естественным сказать: «стек~--- это список, куда добавляют и откуда
забирают только с одного конца»,~--- и написать
\texttt{class Stack<T> extends ArrayList<T>}. Работать это будет, но
\texttt{extends} тянет за собой \emph{все} публичные методы списка: ничто не
помешает вызвать \texttt{stack.add(0, item)} и вставить элемент в середину
стека~--- наследование выставило наружу внутреннюю реализацию и сломало
абстракцию. Композиция решает задачу иначе: список становится не родителем, а
полем (листинг~\ref{lst:03-7}).

\begin{listing}[H]
\begin{minted}{java}
class StackByComposition<T> {
    private final List<T> items = new ArrayList<>();

    public void push(T item) {
        items.add(item);
    }

    public T pop() {   // забираем только с вершины
        if (items.isEmpty()) {
            throw new IllegalStateException("Стек пуст");
        }
        return items.remove(items.size() - 1);
    }
}
\end{minted}
\caption{Стек, построенный на композиции}
\label{lst:03-7}
\end{listing}

Теперь наружу видны ровно те методы, которые нужны стеку. Это первое различие
двух подходов: наследование показывает всем \texttt{public}- и
\texttt{protected}-члены родителя, композиция~--- только объявленное самим
классом. Второе: родитель задан в коде навсегда, а объект в поле подменяют
снаружи~--- \texttt{ArrayList} внутри можно заменить на \texttt{LinkedList}, и
вызывающий код не изменится. Третье~--- \term{проблема хрупкого базового класса}
(\eng{fragile base class problem}): автор родителя меняет внутреннее
устройство, не трогая внешнего контракта, и поведение всех потомков молча
меняется вместе с ним.

Отсюда правило: \emph{предпочитайте композицию наследованию}. Наследуйтесь,
когда отношение is-a действительно есть и оба класса пишете вы; чтобы
переиспользовать чужой код, берите композицию. Не случайно старый
\texttt{java.util.Stack} унаследован от \texttt{Vector} и страдает ровно этой
болезнью.

\section{Запрет наследования: final и sealed}

Иногда наследование нужно, наоборот, ограничить. Простейший инструмент~--- слово
\texttt{final}: переменную нельзя присвоить второй раз, метод нельзя
переопределить, класс нельзя расширить (так объявлены \texttt{String} и
\texttt{Math}).

\begin{oshibka}
\textbf{Частая ошибка: \texttt{final} не делает объект неизменяемым.} Запись
\texttt{final List<String> items = new ArrayList<>();} запрещает присвоить
переменной \emph{другой} список, но вызов \texttt{items.add(...)} по-прежнему
разрешён: неизменна ссылка, а не объект, на который она указывает.
\end{oshibka}

Java~17 добавила вариант между «наследуйся кто хочет» и «наследоваться
нельзя»~--- \term{запечатанные классы} (\eng{sealed classes}), перечисляющие
наследников поимённо после слова \texttt{permits} (листинг~\ref{lst:03-8}).

\begin{listing}[H]
\begin{minted}{java}
sealed class Vehicle permits Car, Bicycle { }

// ветка закрыта: потомков у Car не будет
final class Car extends Vehicle { }

// ветка открыта: у Bicycle потомки разрешены
non-sealed class Bicycle extends Vehicle { }

class MountainBike extends Bicycle { }
\end{minted}
\caption{Запечатанный класс и его наследники}
\label{lst:03-8}
\end{listing}

Каждый наследник обязан объявить, что будет дальше: \texttt{final}~---
наследования больше нет, \texttt{sealed}~--- снова ограниченный список,
\texttt{non-sealed}~--- ограничение снято. Выигрыш не только в дисциплине: раз
компилятор знает все подтипы, \texttt{switch} по запечатанному типу пишут без
ветки \texttt{default}, а при добавлении нового класса компилятор сам укажет,
где его забыли обработать.

\section{Инкапсуляция}

Автомобиль скрывает двигатель за педалями и рулём: клапанами водитель управлять
не может~--- только газом. Это и есть \term{инкапсуляция}: состояние спрятано
модификатором \texttt{private}, а снаружи доступен лишь набор проверенных
операций. Область действия модификаторов приведена в табл.~\ref{tab:03-1}.

\begin{table}[htbp]
\caption{Область действия модификаторов доступа}
\label{tab:03-1}
\centering
\begin{tabular}{L{4.2cm}cccc}
\toprule
Модификатор & Свой класс & Свой пакет & Потомок & Все \\
\midrule
\texttt{private} & да & нет & нет & нет \\
без модификатора & да & да & нет & нет \\
\texttt{protected} & да & да & да & нет \\
\texttt{public} & да & да & да & да \\
\bottomrule
\end{tabular}
\end{table}

Строка «без модификатора»~--- это \term{пакетная видимость}
(\eng{package-private}), которую поле получает, если про доступ забыли: она
шире, чем кажется. А \texttt{protected} открывает доступ потомкам даже в чужих
пакетах.

Смысл инкапсуляции не в том, чтобы к каждому полю приписать пару методов
\texttt{get} и \texttt{set}~--- такая пара не прячет ничего,~--- а в том, чтобы
объект нельзя было привести в противоречивое состояние
(листинг~\ref{lst:03-9}).

\begin{listing}[H]
\begin{minted}{java}
public class BankAccount {
    private double balance;  // менять - только через методы

    public double getBalance() {  // читать разрешено всем
        return balance;
    }

    public void withdraw(double amount) {
        if (amount > balance) {
            throw new IllegalStateException("Мало средств");
        }
        balance -= amount;
    }
}
\end{minted}
\caption{Состояние, защищённое проверками}
\label{lst:03-9}
\end{listing}

Обойти проверку невозможно. Будь поле \texttt{balance} объявлено
\texttt{public}, любая строка любого класса могла бы записать в него минус
миллион, и виновника пришлось бы искать по всему проекту.

\section{Полиморфизм}

Кнопка «Сохранить» работает одинаково для документа, таблицы и презентации:
вызов один, а реализация у каждого типа файла своя. Это \term{полиморфизм}~---
способность одного вызова приводить к разному поведению.

\term{Переопределение} (\eng{overriding})~--- замена реализации метода в
потомке (листинг~\ref{lst:03-10}): переменная объявлена типом \texttt{Shape}, а
вызывается метод того класса, объект которого создан на самом деле.

\begin{listing}[H]
\begin{minted}{java}
import java.util.List;

abstract class Shape {
    public abstract double area();
}

class Circle extends Shape {
    private final double r;
    public Circle(double r) { this.r = r; }
    @Override
    public double area() { return Math.PI * r * r; }
}

class Rect extends Shape {
    private final double w, h;
    public Rect(double w, double h) { this.w = w; this.h = h; }
    @Override
    public double area() { return w * h; }
}

class Demo {
    public static void main(String[] args) {
        List<Shape> shapes = List.of(new Circle(5), new Rect(4, 6));
        // тип переменной - Shape, метод - настоящего объекта
        for (Shape s : shapes) {
            System.out.println(s.area());
        }
    }
}
\end{minted}
\caption{Полиморфный вызов метода}
\label{lst:03-10}
\end{listing}

Метод выбирается по фактическому типу объекта в момент вызова~--- это
\term{позднее связывание} (\eng{late binding}). Цикл ничего не знает о кругах и
прямоугольниках, поэтому новая фигура его не затронет.

Переопределять метод разрешено не как угодно: сигнатура обязана совпасть,
возвращаемый тип может быть только тем же или более узким
(\term{ковариантным}), сужать видимость и добавлять проверяемые исключения
нельзя. Аннотация \texttt{@Override} не обязательна, но пишите её всегда: она
ловит опечатку в имени, которая иначе тихо создала бы новый метод.

\term{Перегрузка} (\eng{overloading})~--- совсем другое: несколько методов с
одним именем и разными списками параметров в одном классе, например
\texttt{print(String)} и \texttt{print(int)}. К наследованию она отношения не
имеет, и нужную версию компилятор выбирает по типам аргументов ещё до запуска.
Поэтому переопределение называют полиморфизмом времени выполнения, а
перегрузку~--- времени компиляции.

Приведение вверх по иерархии, к предку (\eng{upcasting}), выполняется само
собой: собака гарантированно является животным. Приведение
вниз, к потомку (\eng{downcasting}), требует явных скобок и может закончиться
исключением \texttt{ClassCastException}: строка
\texttt{String s = (String) obj;} для \texttt{obj}, хранящей число,
компилируется, но при выполнении рушится. С Java~16 проверка и приведение
слились в одну запись: в блоке \texttt{if (animal instanceof Dog dog)}
переменная \texttt{dog} уже имеет тип \texttt{Dog}. Длинная лестница
\texttt{instanceof} обычно означает, что полиморфизм применён не там: вместо
разбора типов стоило объявить метод, который каждый класс реализует по-своему.

\praktikum{}

Задания выполняются на JDK~21 или новее: \texttt{switch} с сопоставлением по
образцу на более старых версиях не скомпилируется. Каждую программу доводите
до запуска и сохраняйте вывод.

\subsection*{Задание 1. Ветвление и выбор варианта}

Напишите цепочкой \texttt{if} метод
\texttt{static String classify(int n)}, возвращающий одну из строк:
отрицательное, ноль, однозначное, двузначное, трёхзначное, большое (от 1000).
Проверьте на числах $-5$, 0, 7, 42 и 1000. Затем напишите метод
\texttt{static String getGrade(int score)} в двух вариантах: через классический
\texttt{switch} с \texttt{break} и через стрелочный \texttt{switch} как
выражение. Диапазоны: 90--100~--- отлично, 80--89~--- хорошо, 70--79~---
удовлетворительно, остальное~--- неудовлетворительно. Подсказка: чтобы диапазон
стал одним значением \texttt{case}, разделите оценку на 10 нацело.

Напишите метод \texttt{static String describe(Object obj)} на \texttt{switch} с
сопоставлением по образцу, обработав случаи \texttt{null}, \texttt{Integer}
(отдельно большое~--- через \texttt{when}), \texttt{String}, \texttt{Double} и
всё прочее. Проверьте на \texttt{null}, 42, 300, слове, числе 3,14 и значении
\texttt{true}.

\subsection*{Задание 2. Циклы}

Выведите таблицу умножения от 1 до 10 вложенными циклами \texttt{for}, выровняв
числа по столбцам методом \texttt{System.out.printf} (спецификатор
\texttt{\%4d} отводит числу четыре знакоместа). Реализуйте
\texttt{fibWhile(int n)} и \texttt{fibFor(int n)}~--- $n$-е число Фибоначчи
итеративно: напечатайте значения от $F(0)$ до $F(15)$ и найдите первое число
последовательности, превышающее 1000.

Наконец, наберите поиск в матрице из листинга~\ref{lst:03-5}, объявив матрицу
3 на 3, искомое число и переменные \texttt{row} и \texttt{col} со значением
$-1$, и ответьте письменно, что изменится, если убрать метку у \texttt{break}.

\subsection*{Задание 3. Иерархия автопарка}

Абстрактный класс \texttt{Vehicle} хранит закрытые поля \texttt{brand},
\texttt{model} и \texttt{fuelLevel} (от 0 до 1), даёт геттеры, сеттер
\texttt{fuelLevel} с проверкой диапазона, абстрактные методы
\texttt{getFuelConsumption()} (литров на 100~км) и \texttt{getType()} и готовый
метод \texttt{calculateFuelNeeded(double km)}. Класс \texttt{Car} добавляет
поле \texttt{automatic} с расходом 9,5 для автомата и 8,0 для механики, класс
\texttt{Truck}~--- поле \texttt{cargoTons} с расходом
$20 + 3 \cdot \text{cargoTons}$. Класс \texttt{ElectricCar} наследует
\texttt{Car}, хранит закрытые поля \texttt{batteryLevel} (0--100~\%) и
\texttt{maxRangeKm} и реализует интерфейс \texttt{Electric}: метод
\texttt{charge(double hours)} добавляет 20~\% заряда в час, но не более
100~\%, метод \texttt{getRangeKm()} возвращает
\texttt{maxRangeKm * batteryLevel / 100}, расход топлива равен нулю.

Соберите четыре машины в \texttt{List<Vehicle>} и одним циклом напечатайте для
каждой тип и расход топлива на 500~км, а затем, применив \texttt{instanceof} с
образцом, допечатайте запас хода электромобилей.
Объясните письменно, почему этот цикл не пришлось менять из-за
\texttt{ElectricCar}. Наконец, перегрузите \texttt{calculateFuelNeeded}
вариантом \texttt{(double km, double speedFactor)} и покажите, какую версию
компилятор выбирает при вызове с одним и с двумя аргументами.

\subsection*{Задание 4. Инкапсуляция и композиция}

Доведите \texttt{BankAccount} из листинга~\ref{lst:03-9} до рабочего класса.
Добавьте закрытые поля \texttt{pin} (без геттера),
\texttt{failedAttempts} и \texttt{blocked}, а метод
\texttt{withdraw(String pin, double amount)} научите отказывать по
заблокированному счёту, считать неудачные попытки и блокировать счёт после
третьей, а при верном PIN~--- сбрасывать счётчик и проверять сумму. Объясните
письменно, почему открытый геттер PIN свёл бы защиту на нет.

Вторая часть~--- сравнение двух подходов. Напишите с методами \texttt{push} и
\texttt{pop} класс \texttt{StackByInheritance}, расширяющий
\texttt{ArrayList<T>}, и класс \texttt{StackByComposition} из
листинга~\ref{lst:03-7}. Наполните оба тремя числами, вызовите у первого
унаследованный \texttt{add(0, 999)} и покажите, что \texttt{pop()} возвращает
уже не то, что клали последним; у второго такой строки не будет~---
закомментируйте её и припишите текст ошибки компилятора.

\subsection*{Результаты занятия}

По итогам занятия у вас должны быть файлы \texttt{.java} по всем четырём
заданиям с сохранённым выводом каждой программы и письменные ответы на
заданные в них вопросы.

\kontrolvoprosy

\begin{enumerate}
\item Чем стрелочный \texttt{switch} отличается от классического и какую
      частую ошибку он устраняет?
\item Когда \texttt{switch} становится выражением и зачем в нём слово
      \texttt{yield}?
\item В каком порядке выполняются конструкторы при создании объекта потомка и
      почему \texttt{super(...)} обязан стоять первой инструкцией?
\item Чем цикл \texttt{do-while} отличается от \texttt{while}, а
      \texttt{break} с меткой~--- от обычного \texttt{break}?
\item Почему нельзя удалять элементы коллекции внутри цикла \texttt{for-each}
      и как сделать это правильно?
\item Чем композиция отличается от агрегации и от зависимости? Приведите по
      примеру на каждое отношение.
\item В чём состоит проблема хрупкого базового класса и какие методы «утекают»
      наружу, если построить стек через \texttt{extends ArrayList}?
\item Чем \texttt{final} класс отличается от \texttt{sealed} класса и какие
      модификаторы обязан иметь наследник запечатанного класса?
\item Чем переопределение отличается от перегрузки и зачем нужна аннотация
      \texttt{@Override}?
\end{enumerate}

\testzadaniya

\begin{enumerate}

\item Переменная \texttt{day} равна 2, ветки \texttt{case 1},
      \texttt{case 2} и \texttt{case 3} классического \texttt{switch}
      печатают Пн, Вт и Ср, а \texttt{break} стоит только после ветки
      \texttt{case 3}. Что будет напечатано?
  \begin{enumerate}
    \item[а)] Вт
    \item[б)] ПнВт
    \item[в)] ПнВтСр
    \item[г)] ВтСр
  \end{enumerate}

\item Переменная \texttt{Object obj} содержит строку. Какая ветка сработает в
      \texttt{switch} с ветками \texttt{case Integer i},
      \texttt{case String s} и \texttt{default}?
  \begin{enumerate}
    \item[а)] \texttt{default}
    \item[б)] \texttt{case String s}
    \item[в)] \texttt{case Integer i}
    \item[г)] код не скомпилируется
  \end{enumerate}

\item Сколько раз выполнится тело цикла
      \texttt{do \{ x += 10; \} while (x < 10);}, если до цикла \texttt{x}
      равно 10?
  \begin{enumerate}
    \item[а)] один раз
    \item[б)] ни разу
    \item[в)] два раза
    \item[г)] бесконечно
  \end{enumerate}

\item Вложенные циклы по \texttt{i} и \texttt{j} идут от 0 до 2, внешний
      помечен меткой \texttt{outer}; тело внутреннего начинается строкой
      \texttt{if (j == 1) continue outer;}, а иначе печатается пара цифр.
      Что будет напечатано?
  \begin{enumerate}
    \item[а)] 00 01 10 11 20 21
    \item[б)] 00 01 02 10 11 12 20 21 22
    \item[в)] 00 10 20
    \item[г)] 00
  \end{enumerate}

\item Класс \texttt{Dog} наследует \texttt{Animal}; конструктор
      \texttt{Animal} печатает Animal, конструктор \texttt{Dog}~--- Dog. Что
      выведет создание объекта \texttt{Dog}?
  \begin{enumerate}
    \item[а)] Dog
    \item[б)] Animal Dog
    \item[в)] Dog Animal
    \item[г)] ошибка компиляции
  \end{enumerate}

\item Где обязан располагаться вызов \texttt{super()} или \texttt{this()} в
      конструкторе?
  \begin{enumerate}
    \item[а)] в любом месте конструктора
    \item[б)] в последней строке конструктора
    \item[в)] в первой строке конструктора
    \item[г)] после инициализации всех полей
  \end{enumerate}

\item Какие модификаторы обязан иметь наследник \texttt{sealed}-класса?
  \begin{enumerate}
    \item[а)] только \texttt{public} или \texttt{private}
    \item[б)] \texttt{abstract} или \texttt{static}
    \item[в)] \texttt{final}, \texttt{sealed} или \texttt{non-sealed}
    \item[г)] никаких специальных модификаторов не требуется
  \end{enumerate}

\item Закрытое поле \texttt{balance} равно 100, а сеттер присваивает значение
      только при условии \texttt{b >= 0}. Что вернёт геттер после вызова
      сеттера со значением $-50$?
  \begin{enumerate}
    \item[а)] $-50$
    \item[б)] 0
    \item[в)] 100
    \item[г)] ошибка компиляции
  \end{enumerate}

\item Переменная объявлена как \texttt{Animal a = new Cat();}, метод
      \texttt{speak()} переопределён в классе \texttt{Cat}. Метод какого класса
      вызовет строка \texttt{a.speak()}?
  \begin{enumerate}
    \item[а)] \texttt{Animal}~--- по типу переменной
    \item[б)] ни один: код не скомпилируется
    \item[в)] определяется случайно при каждом запуске
    \item[г)] \texttt{Cat}~--- по типу объекта
  \end{enumerate}

\item Переменная \texttt{Object obj} содержит объект \texttt{Integer}. Что
      произойдёт при выполнении строки \texttt{String s = (String) obj;}?
  \begin{enumerate}
    \item[а)] \texttt{ClassCastException} во время выполнения
    \item[б)] переменная \texttt{s} получит значение 42
    \item[в)] ошибка компиляции
    \item[г)] переменная \texttt{s} получит значение \texttt{null}
  \end{enumerate}

\end{enumerate}
